% Lösungen Zone (zusammenbau v0.1)
\einheitenkopf[][Kennst du schon]{Das kennst du schon}

Zuordnung: 1 → Einheit 1 und 2 setzen darauf auf · 2 → Einheit 1 (Tupel zählen), Einheit 2 (Ergebnistabelle zweier Würfel) und Einheit 5 (Reihenfolgen zählen) · 3 → alle Einheiten · 4 → Einheit 1 (deuten), Einheit 6 (Ast als bedingter Anteil) und Einheit 7 · 5 → Einheit 2, 6 und 8 · 6 → Grundlage der Einheit 3, 5, 6 und 7 · 7 → Grundlage der Einheit 4, 5 und 8 · 8 → Voraussetzung der Typen „Term und Ereignis“ in Einheit 7 (Terme mit Binomialkoeffizienten und Bernoulli-Summen lesen) und der Mammutbäume in Einheit 5 (Anzahl der Reihenfolgen, Lotto-Modell) · 9 → Einheit 8 · 10 → Grundlage der Einheit 4, 5 und 8 · 11 → Grundlage der Einheit 4, 5 und 8

\erg{1}{a) \{KK, KZ, ZK, ZZ\} \quad b) 7; 5; 8 (rot; blau; gelb)}
\erg{2}{a) 12 \quad b) $4 \cdot 3 \cdot 2 \cdot 1 = 24$}
\erg{3}{a) $\frac{2}{12} = \frac{1}{6}$ \quad b) $\frac{60}{336} = \frac{5}{28}$}
\erg{4}{a) Die Augenzahl ist ungerade. \quad b) $\frac{0{,}12}{0{,}4} = 0{,}3$}
\erg{5}{a) 10 \quad b) $0{,}15 \cdot 360^\circ = 54^\circ$}
\erg{6}{a) $\frac{2}{5} \cdot \frac{2}{5} = \frac{4}{25}$ \quad b) $0{,}7$ an jedem Ast B; die zweite Stufe wie die erste (A $0{,}3$, B $0{,}7$)}

6b)

\baumzwei{A/0{,}3,B/0{,}7}{A/0{,}3,B/0{,}7}{A/0{,}3,B/0{,}7}

\erg{7}{a) $\frac{5}{8} \cdot \frac{4}{7} = \frac{5}{14}$ \quad b) $\frac{4}{10} \cdot \frac{3}{9} \cdot \frac{2}{8} = \frac{1}{30}$}
\erg{8}{a) $720$ \quad b) „8 über 3“ $= 56$}
\erg{9}{a) $x = 0{,}8$ \quad b) $x_1 = 0{,}4$, $x_2 = 0{,}6$}
\erg{10}{a) Lisa rechnet mit Zurücklegen; der zweite Zug hat nur noch 3 grüne von 9: $\frac{4}{10} \cdot \frac{3}{9} = \frac{2}{15}$}
\erg{11}{a) $\frac{3}{8} \cdot \frac{2}{7} = \frac{3}{28}$}
